% Imporditud: tools/import_eksperimendid.py
% Allikas: Eesti füüsikaolümpiaadi eksperimendi ülesannete
%   kogu (Rael Kalda, 2023), ülesanne Ü24, lahendus L24.
\setAuthor{Erkki Tempel}
\setRound{piirkonnavoor}
\setYear{2016}
\setNumber{P E2}
\setDifficulty{1}
\setTopic{varia}
\setVahendid{20-sendine euromünt (m = 5,7 g), 5-sendine euromünt, kolmnurkne joonlaud}
\setPunktid{12}
\setAllikas{Eksperimendi kogu Ü24, lahendus lk 42}
\setLahendusseis{toores}

\prob{Pangaröövel}
Pangaröövel varastas 10 kilogrammi 5-sendiseid euromünte. Leidke sellise rahakoguse väärtus eurodes.

\hint

\solu
Selleks, et leida 10 kg 5-sendiste müntide arv, peame leidma ühe mündi massi. Selleks kasutame kangimeetodit. Kuna joonlaud on kolmnurkne, peame kõigepealt leidma selle raskuskeskme. Selleks asetame joonlaua laua servale (joonlaua skaala on laua servaga risti) ning leiame koha, kus joonlaua laua peal olev osa ning üle laua olev osa on tasakaalus (1 p). Edaspidistes mõõtmistes peab joonlaua asend olema laual nii, et joonlaua raskuskeskme asukoht ühtiks laua servaga ning joonlaua skaala on laua servaga risti (2 p). Asetame ühe mündi (näiteks 20-sendise) joonlaua laua peal oleva osa peale ning teise (5-sendise) mündi joonlaua teise otsa peale. Nüüd hakkame liigutama 20sendist joonlaua keskpunkti poole ning leiame koha, kus joonlaud jääb tasakaalu (2 p). Vastavalt kangi reeglile on mündile mõjuvate raskusjõudude ja jõuõlgade korrutised võrdsed.

m20g $\cdot$ l20 = m5g $\cdot$ $l_5$ (1 p)

Mõõtes müntide keskpunktide kaugused (l20 ja $l_5$) pöörlemisteljest (2 p) ning teades 20-sendise mündi massi, saame 5-sendise mündi massiks

$m_5$ = 5,7 g $\cdot$ l20 (1 p) $l_5$

Vastus täpsusega 3,9 g $\pm$ 0,2 g annab 2 p, 3,9 g $\pm$ 0,4 g annab 1 p. Seega 10 kg müntide väärtus on

n= 10 kg $\cdot$ 0,05 EUR $\approx$ 130 EUR. (1 p)

$m_5$
\probend
